\documentclass{ltjsarticle}
\usepackage[utf8]{inputenc}
\usepackage{ascmac}
\usepackage{amsmath, amssymb}
\usepackage{amsfonts}
\usepackage{amsthm}
\usepackage{url}
\usepackage{hyperref}
\theoremstyle{definition}
\newtheorem{thm}{定理}
\newtheorem{dfn}{定義}
\newtheorem{prob}{練習問題}
\newtheorem{expr}{例題}
\makeatletter
\renewenvironment{proof}[1][\proofname]{\par
  \pushQED{\qed}%
  \normalfont \topsep6\p@\@plus6\p@\relax
  \trivlist
  \item\relax
  {
  #1\@addpunct{.}}\hspace\labelsep\ignorespaces
}{%
  \popQED\endtrivlist\@endpefalse
}
\makeatother
\usepackage[dvipdfmx]{graphicx}
\title{Prime Constraint}
\author{じゅんにー}
\begin{document}
\maketitle
\tableofcontents
\newpage
\section{前提知識}
知っていると読みやすくなる程度にみればヨシ！(知らないなら出てきたときに調べれば良いだろう)

\begin{thm}(フェルマーの小定理)\\
$p$を素数,$a$を$p$と互いに素な整数とすると$$a^{p-1}\equiv1\pmod p$$が成立する.
\end{thm}

\begin{thm}(オイラーの定理)\\
$m$を整数,$a$を$m$と互いに素な整数, $\varphi(m)$で$m$以下の$m$と互いに素な正整数の個数を表すものとすると$$a^{\varphi(m)}=1\pmod m$$が成立する.
\end{thm}

\begin{dfn}($p$進付値)\\
$p$を素数とし,$n$を整数とするとき$v_p(n)$で$n$が$p$で割り切れる回数を表す.(便宜的に$v_p(0)=\infty$と定める)
\end{dfn}

\begin{thm}(LTEの補題)\\
$p$を素数$a,b$を$p$で割り切れない整数とし, $v_p(a-b)>\dfrac1{p-1}$が成立するとする. このとき任意の非負整数$n$について$$v_p(a^n-b^n)=v_p(a-b)+v_p(n)$$が成立する.
\end{thm}

\begin{dfn}(平方剰余記号/ルジャンドル記号)\\
$p$を奇素数, $a$を整数とする.\\
(1)\ $a\equiv0\pmod p$なら$\left(\dfrac ap\right)=0$\\
(2)\ $a\not\equiv0\pmod p$かつ$a\equiv x^2\pmod p$なる整数$x$が存在するなら$\left(\dfrac ap\right)=1$\\
(3)\ $a\not\equiv0\pmod p$かつ$a\equiv x^2\pmod p$なる整数$x$が存在しないなら$\left(\dfrac ap\right)=-1$\\
と定義する.
\end{dfn}

以下は証明を知らなくても良いだろう.

\begin{thm}(平方剰余の諸法則)\\
$p,q$を奇素数とするとき\\
(1)\ $\left(\dfrac{-1}p\right)=(-1)^{\frac{p-1}2}$\ (第一補充則)\\
(2)\ $\left(\dfrac{2}p\right)=(-1)^{\frac{p^2-1}8}$\ (第二補充則)\\
(3)\ $\left(\dfrac qp\right)\left(\dfrac pq\right)=(-1)^{\frac{p-1}2\frac{q-1}2}$\ (相互法則)
\end{thm}

あと知っておくと良いのは原始根の存在定理ぐらいだろう.

\begin{thm}(原始根の存在)\\
$p$を素数とするとある整数$a$が存在して任意の$p$で割り切れない整数$x$に対しある整数$m$が存在して$a^m\equiv x\pmod p$となる.
\end{thm}

\section{概論}
ここでは比較的高難易度な不定方程式に頻出の技法であるprime-constraintについて解説する.なお,これは筆者の造語なのでここ以外では通じないと思う.この手法は適当な素因数(もしくは約数等)に関する条件によって不定方程式の解を絞る方法である.まずは簡単な例を見る.

\begin{expr}(IMOLonglist1992P63改題)\\
(1)\ $b$を正の奇数とするとき$b^2-50$が$4$で割って$3$余る素因数を持つ条件を求めよ.\\
(2)\ $\displaystyle\frac{2a^2+1}{b^2-50}$が整数となるような非負整数の組$a,b$を全て求めよ.
\end{expr}

(1)で適当な素因数の存在条件を見て, (2)でそれを不定方程式に適用しているのが分かるだろう. (1)では$b^2-50\equiv3\pmod4$である事が使えそう, (2)は平方剰余を使えば良さそうである.\\
\textbf{解答}\\
(1)\ $b=1,3,5,7$のときは調べると$b=1$のみが条件を満たす. $b>7$のとき$b^2-50>0$であり, $b$は奇数なので$b^2-50\equiv3\pmod4$である. $b^2-50$が$4$で割って$3$余る素因数を持たないと仮定すると素因数分解を考えて$$3\equiv b^2-50=\prod_{pは素数,\ p|b^2-50}p^{v_p(b^2-50)}\equiv1\pmod4$$これは矛盾である. よって$b=1,b\geq9$においては$4$で割って$3$余る素因数を持つ.\\
($b^2-50<0$のときは素因数分解の先頭に$-$が付くので上の議論は成立しないことに注意)\\
(2)\ $2a^2+1$は奇数なので$b^2-50$も奇数, 特に$b$は奇数である. $b^2-50$が$4$で割って$3$余る素因数$p$を持つと仮定すると$b^2-50\equiv0\pmod p,\ 2a^2+1\equiv0\pmod p$が成立する. よって$(2ab)^2\equiv 2\cdot b^2\cdot 2a^2\equiv 2\cdot 50\cdot (-1)\equiv-100\pmod p$となる. しかし$p\equiv3\pmod4$より$p\neq2,5$にも注意するとこれは平方剰余の第一補充則に矛盾する. よって$b^2-50$は$4$で割って$3$余る素因数を持たない. つまり(1)より$b=3,5,7$でありこれらを平方剰余の第一第二補充則を用いながら条件を満たす$a$を探すと$(a,b)=(41k+15,3),(41k+26,3),(k,7)$($k$は非負整数)が求める解である.\\

この問題から分かるようにこの手法は,\\
\textbf{1.}この手法で解くことが出来ることに気付く\\
\textbf{2.}満たして欲しい条件を考え,素因数の存在を言う\\
\textbf{3.}実際に問題に適用して解を絞る\\
の大体$3$パートに分かれている. そして, 知っておくべきはどのような問題にどのような条件を使うと上手くいくことが多いのかである. 以下では典型的な$5$つの例, フェルマーの小定理との組み合わせ, 素因数(約数)の最小性の利用, 平方剰余の利用,無限に大きく素数を取る方法, ジグモンディーの定理の利用の問題について紹介する.(ジグモンディーの定理の利用は知識ゲー感が強いので数オリに頻出というわけではないが) なお, 冒頭の例題は三つ目の平方剰余の利用に該当する. (というのも他はより考察に前提知識を要求するためこの手法の説明には向かないので.)

\section{フェルマーの小定理との組み合わせ}
フェルマーの小定理との組み合わせと言われただけで何をすべきか分からないと思うのでまずは有名問題で手法を見ていこう.

\begin{expr}\ \\
$p$を素数とする.\\
(1)\ $n$を$1$ではない整数とすると$\displaystyle\frac{n^p-1}{n-1}$の素因数は$p$であるか$p$で割って$1$余る事を示せ.\\
(2)\ $p$で割って$1$余る素数が無限に存在することを示せ.
\end{expr}

(2)はディリクレの算術級数定理の一部である.(多分別の場所で解説するので詳細は省略)有名な素数の無限性の証明をなぞれば証明できる.\\
\textbf{解答}\\
(1)\ $q$を$\displaystyle\frac{n^p-1}{n-1}$の$p$でなく, $p$で割った余りが$1$ではない素因数とする. このとき$n^p\equiv1\pmod q$が成立する. また, これより$n$は$q$で割り切れないのでフェルマーの小定理より$n^{q-1}\equiv1\pmod q$も成立する. 仮定より$(q-1,p)=1$なので$kp\equiv 1\pmod{q-1}$なる正整数$k$が存在し, $1\equiv n^{kp}\equiv n\pmod q$が成立する. 特に$n\equiv 1 \pmod q$である. しかしこれは$0\equiv\displaystyle\frac{n^p-1}{n-1}=n^{p-1}+...+n+1\equiv1^{p-1}+...+1+1=p\pmod q$となるので矛盾.(ここにおいてはLTEの補題を用いても良い)よって題意は示された.\qed\\
(2) $p$で割って$1$余る素数が有限個だと仮定し,その総積を$S$とする.(1)で$n=pS$とすると$\displaystyle\frac{n^p-1}{n-1}$は$p$で割り切れず, また$S$の取り方により$p$で割って$1$余る素因数も持たない. しかし(1)よりこれは$\displaystyle\frac{n^p-1}{n-1}=\pm1$である事を意味するがこれは明らかに不合理である. よって$p$で割って$1$余る素数は無限に存在する.\qed\\

この証明では指数が素数であって約数のバリエーションが少ないことが鍵である. というのも$n\equiv1\pmod q$を示すのに用いた事実はより一般に次のようになるためである.

\begin{thm}\ \\
$x,a,b,m$を正整数とし$x^a\equiv x^b\equiv 1\pmod m$が成り立つなら$x^{(a,b)}\equiv1\pmod m$が成立する.
\end{thm}

これは上の証明と同様に$ak-bl=1$なる非負整数$a,b$を考えれば示せる.($(\mathbb Z/m\mathbb Z)^*$の群構造を考えればより見通しよく定理を見られるだろう)\\
ここで最大公約数が出てくるので約数が少ない方が得なのである. なお, 次の類題から分かるよう, 約数が多いときでも円分多項式に近ければ上手くいく.

\begin{prob}\ \\
$\phi_n(x)$を$n$次円分多項式とする.\\
$n$を正整数, $x$を整数とする. $\phi_n(x)$の$n$と互いに素な素因数を$n$で割った余りは$1$である事を示せ.
\end{prob}

\begin{prob}(ISL2006N5, 2007春P2)\\
$\displaystyle\frac{x^7-1}{x-1}=y^5-1$が成立する$1$でない整数の組$(x,y)$を全て求めよ.
\end{prob}

なお次の例題からも分かるようにこの議論は法が素数でなくても上手くいくこともある. 
\begin{expr}(APMO2012P3)\\
$\displaystyle\frac{n^p+1}{p^n+1}$が整数となるような素数$p$と正整数$n$の組を全て求めよ.
\end{expr}

分子の指数が素数なのでprime-constraintが使えそうだが, その前に大小関係や偶奇などの基本情報は把握しておくべきである.\\
\textbf{解答}\\
$p=2$のとき$n=2,4$が条件を満たす. $p\geq3$のとき分母は偶数なので分子も偶数. よって$n$は奇数である. 明らかに$n\neq1$より$n\geq3$. $\displaystyle\frac{n^p+1}{p^n+1}\geq1$より$n^p+1\geq p^n+1$. これは$n\geq3,\ p\geq3$より(対数を取って微分すれば分かるが)$p\geq n$を意味する.

ここで上記の手法を用いる. $n$は奇数なので分母は$p+1$で割り切れる. よって$n^p+1\equiv0\pmod{p+1}$. 特に$n^{2p}\equiv1\pmod{p+1}$.  またオイラーの定理より$n^{\varphi(p+1)}\equiv1\pmod{p+1}$なので$n^{(2p,\varphi(p+1))}\equiv1\pmod{p+1}$. $p\geq3$なので$\displaystyle\varphi(p+1)\leq\frac{p+1}2<p$より$(2p,\varphi(p+1))=1,2$. $\varphi(p+1)$は偶数なので$(2p,\varphi(p+1))=2$で, よって$n^2\equiv1\pmod{p+1}$. よって$0\equiv n^p+1\equiv n+1\pmod{p+1}$. 特に$n+1\geq p+1$. 最初に得られた$p\geq n$と合わせて$n=p$. このとき条件を満たす.

よって求める解は$(n,p)=(p,p),(4,2)$である.\\

いくつかコメントしておく. まずは$n^p+1\equiv0\pmod{p+1}$から$n^{2p}\equiv1\pmod{p+1}$を導いて,その後にオイラーの定理を適用した部分, これは$a^p+1$型に対してprime-constraintを行う際良くやる定石である. この議論を一般化すれば以下のようになる.

\begin{thm}\ \\
$x,a,b,m$を正整数とし, $x^a\equiv-1\pmod m,\ x^b\equiv1\pmod m$が成り立っているとする. \\このとき$x^{(a,b)}=-1\pmod m$が成立する. また$v_2(b)\leq v_2(a)$ なら$m=1,2$となる.
\end{thm}

また, なぜ素因数ではない約数を用いたのかは, 今回欲しかった性質が$n+1\geq p+1$という不等式だったからと言う部分が大きいと思う.
例えば平方剰余の性質などもっと素数特有の性質を利用したいなら素因数を持ち出すべきである. 今回の場合は, 素因数から言える情報はあまり役に立たず(手を動かせば分かる). 唯一使えたのが大小関係だったという感じである. また, 素因数に制約が付かない場合は素因数を全て掛け合わしても上手くいく事も記しておく.

\begin{prob}(Prothの判定法)\\
$N=h\cdot2^n+1,\ h<2^n$となる正整数$h,n$が存在するとする. このとき$N$が素数である事の必要十分条件は$$a^{\frac{N-1}2}\equiv-1\pmod N$$となる整数$a$が存在することである事を示せ.
\end{prob}

より一般に次が知られている.

\begin{prob}(Pocklginton,\ 1914)\\
$N=FR+1,\ (F,R)=1$かつ$F$の各素因数$q$に対して$a^{N-1}\equiv1\pmod N$かつ$(a^{\frac{N-1}q}-1,N)=1$となる整数$a=a(q)$が存在するとき, $N$の素因数は$F$を法として$1$と合同である. 特に$F+1>\sqrt N$となるとき$N$は素数である.
\end{prob}

\begin{prob}(TST2003, by Reid Barton)\\
素数の組$(p,q,r)$であって$$p|q^r+1,\ q|r^p+1,\ r|p^q+1$$を満たすものを全て求めよ.
\end{prob}

最後に分母の特徴を用いて条件を満たす素因数を見つける方法を見る.

\begin{expr}\ \\
$\displaystyle\frac{n^p-1}{p^2-n^2}$が\textbf{非負}整数となるような素数$p$と正整数$n$を全て求めよ.
\end{expr}

まず明らかに$p>n$である. 分子の素因数は$n-1$の約数か$p$か$p$で割って$1$余る. それに対して分母の約数を$p$で割った余りはほとんど$p$で割って$1$余ることはない. これに気が付けば解けるだろう.\\
\textbf{解答}\\
$n=1$のときは明らかに条件を満たす. また$\displaystyle\frac{n^p-1}{p^2-n^2}\geq0$より$p>n$であるので$1<n<p$とすると$p^2-n^2=(p+n)(p-n)$の素因数を$p$で割った余りは$0,1$にはならない. 
しかし$\displaystyle\frac{n^p-1}{n-1}$の素因数を$p$で割った余りは$0$か$1$である.よって$\displaystyle(p^2-n^2,\frac{n^p-1}{n-1})=1$. よって$\displaystyle\frac{n-1}{p^2-n^2}$は正整数. 従って$n-1\geq p^2-n^2$. しかしこれは$n\leq p-1$なので不適.

よって求める解は$(n,p)=(1,p)$である.\\

(これ実戦で見たことないような...(と思ったら実際に出題された例がありました！)とりあえず自作問題の中から練習問題を出しておきます)

\begin{prob}\ \\
$\cfrac{5^{3pq}-1}{(p^2+10)(3^q-8)}$が整数となるような素数の組$p,q$を全て求めよ.
\end{prob}

\section{最小性の利用}
これはほとんど$1$個の手法に集約されるので後から一個前の節に入れようかと思ったが, 少しだけ前節とは違う部分があるので分けてみた.例題はいわゆるマスターデーモンと呼ばれる問題である.

\begin{expr}(IMO1990P3)\\
$\cfrac{2^n+1}{n^2}$が整数となるような$1$より大きい整数$n$を全て求めよ.
\end{expr}

鍵は$n$の最小の素因数を使うことである.\\
\textbf{解答}\\
分子は奇数より$n$は奇数である. $n>1$より$n$は素因数を持ち, $n$の最小の素因数を$p$とすると$2^{n}\equiv-1\pmod p$. またフェルマーの小定理により$2^{p-1}\equiv1\pmod p$である. よって前節の結果により, $2^{(n,p-1)}\equiv-1\pmod p$. さらに$p$は$n$の最小の素因数なので$(n,p-1)=1$. よって$3\equiv0\pmod p$, つまり$p=3$である.

$n=3^a m$($a$は正整数, $m$は$3$で割り切れない整数)と表すと,分母は$3$で$2a$回割り切れるが,分子はLTEの補題により($n$は奇数である事に注意)$3$で$a+1$回しか割り切れない. よって$2a\leq a+1$から$a=1$を得る. 

$m$が素因数を持つとし, 最小の素因数を$q$とすると, 同様に$2^{(3m,q-1)}\equiv -1\pmod q$. $q$の素因数の最小性と$(-1)^3=-1$に注意して$2^3\equiv -1\pmod q$, つまり$q|9$を得るがこれは$m$が$3$で割り切れないことに矛盾. よって$m=1$となり, 求める整数$n$は$n=3$のみである.\\

指数に素数が来ていなくても最小性などを駆使すればうまくいく事もあるという感じである.ちなみに似たような問題が別の年のIMOに出ている.\\

\begin{prob}(IMO2000P5)\\
丁度$2000$個の相異なる素数で割りきれるような自然数$n$であって, $n$が$2^n+1$を割り切るようなものは存在するか.
\end{prob}

この問題は\cite{精選}にある.
\begin{prob}\ \\
$\dfrac{3^n+1}{n}$が整数となる$3$以上の奇数は存在しないことを示せ.
\end{prob}

\begin{prob}(Russia2000)\\
どの$2$つも互いに素であるような$3$整数$a,b,c>1$であって$$b|2^a+1,\ c|2^b+1,\ a|2^c+1$$を満たすものが存在するかどうかを決定せよ.
\end{prob}

上の問題は底を変えたときに結果がどのように変わるのか研究してみるのも一興だろう.

\begin{prob}(Eurosky 杯P7)\\
$\dfrac{7^n+1}{\varphi(n)}$が整数となるような正整数$n$を全て求めよ.
\end{prob}

上のものとは少し違うがこれも似たカテゴリーに属するだろう.

\begin{prob}(AMO(旭丘数学オリンピック)P13)\\
$\cfrac{p^p+1}n,\ \cfrac{n^n+1}p$がともに整数となるような正の整数$n$と素数$p$の組を全て求めよ.
\end{prob}

この解法はよくある不定方程式だけでなく, 約数の集合に関する問題などにも応用できる.
\begin{expr}(IMO2002P4誘導削除)\\
$n$を$2$以上の整数とする. $n$の正の約数を小さい順に$1=d_1,d_2,...,d_k=n$とするとき$n^2$は$D=d_1d_2+d_2d_3+...+d_{k-1}d_k$で割り切れた. $n$としてあり得る値を全て求めよ.
\end{expr}
\textbf{解答}\\
$n\geq2$より$n$は素因数を持ち, その中で最小のものを$p$とすると$d_{k-1}=\dfrac np$となるから$D\geq d_{k-1}d_k=\dfrac{n^2}p$. 一方
\begin{align*}
    D=&n^2(\dfrac1{d_1d_2}+\dfrac1{d_2d_3}+...+\dfrac1{d_{k-1}d_k})\\
    <&n^2(\dfrac1{1\cdot2}+\dfrac1{2\cdot3}+...+\dfrac1{i(i+1)}+...)\\
    =&n^2((\dfrac11-\dfrac12)+(\dfrac12-\dfrac13)+...+(\dfrac1{i}-\dfrac1{i+1})+...)\\
    =&n^2
\end{align*}
より$D<n^2$. よって$n^2$の$n^2$未満の最大の約数は$\dfrac{n^2}p$である事に注意すれば$D=\dfrac{n^2}p$となるが, これが達成されるのは$n=p$の時のみである. よって条件を満たす$n$は$n=p$\ ($p$は素数)である.\\

またこの方法は構成ゲーでも便利である.
\begin{prob}(2016 Turkey EGMO TST P6)\\
$n$を$2$以上の整数で平方因子を持たないものとする. このとき$n$の素因数$p$と整数$m$であって$n\ |\ p^2+p\cdot m^p$なるものが存在する事を示せ.
\end{prob}

\section{平方剰余の利用}
この節では, 平方剰余を利用するものを見ていく.

\setcounter{expr}{0}
\begin{expr}(再掲)(IMOLonglist1992P63改題)\\
(1)\ $b$を正の奇数とするとき$b^2-50$が$4$で割って$3$余る素因数を持つ条件を求めよ.\\
(2)\ $\displaystyle\frac{2a^2+1}{b^2-50}$が整数となるような非負整数の組$a,b$を全て求めよ.
\end{expr}

\setcounter{expr}{6}
このように$\bmod 3,4$での値を考える$\rightarrow\ \bmod 3,4$で不都合な素因数の存在を言う$\rightarrow$平方剰余の諸法則に矛盾させる, という手法をとることが多い.もう少し問題を見ていこう.

\begin{expr}(2003春P8)\\
次の式を満たす整数$d,m,n,k$の組を全て決定せよ.$$(6d-1)^m+(3d+4)^n=k^2$$
\end{expr}

上の定石通りにやれば解けるが, 場合分けがかなり大変である.\\
\textbf{解答}

まず$d=0$のとき$(-1)^m=(k+2^n)(k-2^n)$となるので$n=0,k=0,m=2t+1$($t$は整数)が条件を満たす.\\
$d=-1$のとき右辺は非負整数より$m$は非負偶数. よって$1=(k+7^{\frac m2})(k-7^{\frac m2})$となり, 解は無い. 以下$d\neq 0,-1$とする.

$(6d-1,3d+4)=(-9,3d+4)=1$より$6d-1,3d+4$は互いに素である. 従って$m,n$が負になり得るのは$d=0,-1$のときなので$m,n$は非負整数であるとして良い.

両辺の$\bmod3$を取ると$(-1)^m+1\equiv0,1\pmod3$となるので$m$は奇数である.\\
$(i)\ d$が偶数のとき

両辺の$\bmod 4$を取って$-1+(3d)^n\equiv0,1\pmod4$. よって$n=0,1$である.

$n=0$のとき, $(6d-1)^m=(k+1)(k-1)$となり, $d>0,\ (k+1,k-1)=1,\ k\equiv0\pmod6$である事から, ある非負整数$e,f$があって$k+1=(6e+1)^m,k-1=(6f-1)^m$と表せる.\ 二式の差をとって$2=(6e+1)^m-(6f-1)^m$. 右辺の因数分解を考えると$m=1$. このとき$d=6t^2,\ k=\pm6t$($t$は非負整数)が条件を満たす.

$n=1$のとき, $d>0$であり, $\bmod8$で考えると$d\equiv6\pmod8$である. よって $3d+4=2\left(3\cdot\dfrac{d}2+2\right)$には$3\cdot\dfrac d2+2\equiv3\pmod4$なので$4$で割って$3$余る素因数$p$が存在する. $6d-1=2(3d+4)-9$より, 与式の両辺の$\bmod p$を取ると$(-9)^m\equiv k^2\pmod p$. これは$-1$が$\bmod p$の平方剰余である事を意味するが, $p\equiv3\pmod 4$と平方剰余の第一補充則に矛盾.

$(ii)\ d$が奇数のとき\\
$6d-1\equiv1\pmod4,\ 3d+4\equiv 3\pmod4$より$n$は奇数. 与式の$\bmod 4$を取ると$9d+3\equiv0\pmod4$. よって$d\equiv1\pmod4$. 明らかに$d>0$なので$3d+4(\equiv3\pmod4)$には$4$で割って$3$余る素因数$p$が存在する事となり$(i)$と同様, 平方剰余の第一補充則に矛盾する.\\
以上より求める解は$(d,m,n,k)=(0,2t+1,0,0),(6t^2,1,0,6t)$($t$は整数)\\

かなり重たかったが, このように議論の途中で平方剰余の第一補充則が刺さるという感じで解ける問題もしばしばある.

\begin{prob}\ \\
$3^k(m^2+k^2)=n!-3$を満たす非負整数$k,m,n$を全て求めよ.
\end{prob}

\begin{prob}\ \\
以下の式を満たす非負整数$a,b,c$を全て求めよ.$$2^a+37^a-3^b=c^2$$
\end{prob}

最後に平方剰余の相互法則を用いる問題を書いておこう.

\begin{expr}(RomaniaTST2019,2020第2回通信P5)\\
$m^6=n^{n+1}+n-1$をみたす正の整数の組$(m,n)$を全て求めよ.
\end{expr}

平方数や立方数で挟める所は挟んで解を絞ることが出来るがそこからprime-constraintを使う部分が難所である.\\
\textbf{解答}

$n$が奇数のとき$(n^{\frac{n+1}2})^2\leq n^{n+1}+n-1<(n^{\frac{n+1}2}+1)^2$より左辺の等号が成立し, $n=1,m=1$.

$n$が偶数の時, $\bmod3$より$n$は$3$の倍数でない. $n\equiv 2\pmod3$のときは$(n^{\frac{n+1}3})^3< n^{n+1}+n-1<(n^{\frac{n+1}3}+1)^3$より不適. 残るは$n\equiv 4\pmod6$のときのみである.$n+1\equiv5\pmod6$より$p|n+1,\ p\equiv2\pmod3,\ p\neq2$なる素数$p$が存在する. この$p$に対して$m^6\equiv-1^{n+1}-1-1\equiv-3\pmod p$. 特に$\displaystyle\left(\frac{-3}p\right)=1$. しかし, 平方剰余の相互法則より$\displaystyle\left(\frac{-3}p\right)=(-1)^{\frac{p-1}2+\frac{p-1}2}\left(\frac p3\right)=\left(\frac 23\right)=-1$となるため, これは不適である.

よって求める解は$(m,n)=(1,1)$のみである.\\

$n+1$の素因数を利用する部分はかなり難しいと思うが$n$の値によらずに値が一定になるような約数の取り方をすると良いという教訓は得られるだろう.

\section{無限に大きく素数を取る方法}(いつかちゃんと書く)
書くこと: ディリクレの算術級数定理, 小林の定理あたりの話, e,piの超越性, 関数方程式, 局所大域寄りの話？

ここまでは主に不定方程式に特化した手法が多かったが, ここでは幅広く使われる素数の無限性を用いる手法について解説する. 素数は$\bmod$や約数の議論がしやすいだけでなく, 沢山存在するという強みもある. まずは素数の見つけ方の一例をその考え方と共に見ていく.

\begin{expr}(ISL2012N6)\\
$x,y$を正整数とする. 任意の整数$n$に対して$x^{2^n}-1$が$2^ny+1$で割り切れるなら, $x=1$である事を示せ.
\end{expr}

まず$x^{2^n}$の指数の約数は二冪に限られているから$2^ny+1$のいい感じな素因数$p$を持ってくれば良さそうだと分かる. 必要な$p$の条件は$x^{2^n}-1\equiv x^{(2^n,p-1)}-1\pmod p$に注意すれば$p-1$があまり$2$で割り切れない, さらに$x$に対して十分大きい事である.

一つ目の条件を考えるとできれば$n$が小さいときの$2^ny+1$の素因数$p$を取りたい($v_2(p-1)$が小さく抑えやすいので). 一方二つ目の条件を見ると$p$ は$n$が大きい時のものをとらないといけない.

そこで実は$n$が大きくても$2^ny+1$の素因数$p$で$v_2(p-1)$が小さいものがあることを示したい. ($p$が小さいと$x$をめっちゃ大きく取ると議論が破綻するので前半の願望には譲歩してもらった)

ここで$n$が大きくても$2^n\pmod p$は周期的だから$p|2y+1$かつ十分大きいある整数$n$に対して$p|2^ny+1$, みたいな感じに出来る. ここで$v_2(\dfrac{2^ny+1}p-1)$は小さいので$\dfrac{2^ny+1}p$の素因数$q$で$v_2(q-1)$が小さいものが存在する. さらに$n$を工夫して取れば$q$は大きく出来る. これはフェルマーの小定理を使って小さい素数$r$に対して$2^ny+1\equiv 2y+1\pmod r$となるように$n$を取ってしまえば良いのである. これらの考察から以下の解答が得られる.(かなり雑ですが)\\
\textbf{解答}\\
$n=1+\varphi((2y+1)^2\cdot x^{2y+1}!)$とすると$2^ny+1\equiv0\pmod{2y+1}$. $2^ny+1\equiv2y+1\pmod{2y+1}$より$\dfrac{2^ny+1}{2y+1}$は整数で$2y+1$で割り切れない. さらに$2y+1$を割り切らない$x^{2y+1}$以下の任意の素数$q$について$2^ny+1\equiv2y+1\not\equiv0\pmod q$. よって$\dfrac{2^ny+1}{2y+1}$の素因数は$x^{2y+1}$より大きい. さらに$\dfrac{2^ny+1}{2y+1}-1\equiv \dfrac{2y}{2y+1}\not\equiv0\pmod{2^{v_2(2y)+1}}$ だから$\dfrac{2^ny+1}{2y+1}$の素因数$p$で$v_2(p-1)\leq v_2(2y)$を満たすものが存在する. 仮定より$x^{2^n}\equiv1\pmod p$なので$x^{2^{v_2(2y)}}-1\equiv0\pmod p$. よって$x^{2y}-1\equiv0\pmod p$. しかし$p>x^{2y+1}$なのでこれは$x=1$を意味する. よって題意は示された.\qed

このように$p$として無限に大きく取れるようなものを使うことで不等式評価を作り, 矛盾を導くことが出来る. $2$から$5$節の技法はここでも必要になってくることにも注意しよう. ここでの解法のイメージには次の原則がある.

\begin{center}
    \textbf{大抵の関数は十分大きいところで十分大きい素因数を持つ}
\end{center}

\begin{prob}(BalkanMO Shortlist)\\
$x,y$を正整数とする. 任意の正整数$n$に対して
$$(ny)^2+1\mid x^{\varphi(n)}-1$$
となるなら$x=1$である事を示せ.
\end{prob}

上の原則は多項式においても顕著である.
\begin{expr}(Ibero America 2019, 2020第1回通信添削P4)

$a_1,a_2,...,a_{2020}$を正の整数とする. $P(x)$を整数係数多項式であって, 任意の正の整数$n$について$P(n)$が$\displaystyle \sum_{i=1}^{2020}a_i^n$の約数となるようなものとする. このとき$P(n)$は定数関数である事を示せ.
\end{expr}

純粋な大小比較では上手く行かなさそうなので合同式関係の議論が必要そうである. $p|P(n)$なる素数$p$を取ってみよう. すると$p|\displaystyle \sum_{i=1}^{2020}a_i^n$だがこれだけでは情報が少なすぎる. この$p$についてもう少しいい話が無いかと探ると$p|P(n+p)$である事に気付くだろう. よって$p|\displaystyle \sum_{i=1}^{2020}a_i^{n+p}$. フェルマーの小定理を使えば$p|\displaystyle \sum_{i=1}^{2020}a_i^{n+1}$. これを繰り返せば$p|\displaystyle \sum_{i=1}^{2020}a_i^{k(p-1)+1}$, つまり$p|\displaystyle \sum_{i=1}^{2020}a_i$と出来る. ($p,a_i$が互いに素である事を仮定していないのでその関係で式が微妙に汚い)ここで$p$が無限に大きく取れればこれは矛盾である.

\textbf{解答}

$P$は定数でない多項式であると仮定して矛盾を導く. まず, 一般に次の定理が知られている. これは初等的に示せる上, たまに使うため定理の形で記しておく.
\begin{dfn}

$X$を整数の部分集合とするとき, $\mathbb P(X)$を$X$のある元を割りきる素数全体の集合とする.

\end{dfn}
\begin{thm}(シューアの定理)

任意の定数でない整数係数多項式$P(x)$に対し$\mathbb P(P(\mathbb N))$は無限集合である.

ここで$P(\mathbb N)=\{P(n)|n\in \mathbb N\}$である.
\end{thm}
\textbf{証明}(素数の無限性の証明のアナロジーである)

$\mathbb P(P(\mathbb N))$が有限集合であると仮定し, その総積を$S$とする. $P$は非有界より$S>1$である.
まず$P$は定数でないので$|P(k)|>0$なる$k\in \mathbb N$が存在する.  $P(k+S^{|P(k)|}\cdot n)\equiv P(k)\pmod{S^{|P(k)|}}$. ここで$P(k+S^{|P(k)|}\cdot n)$の素因数$p$は$S$を割り切り, $v_p(P(k+S^{|P(k)|}\cdot n))=v_p(P(k)+m\cdot S^{|P(k)|})=v_p(P(k))$なので$|P(k+S^{|P(k)|}\cdot n)|\leq |P(k)|$(実際には等号)これは$n$を動かすことを考えると$P$の非有界性に矛盾. よって定理は示された.\qed

定理により$p|P(n)$かつ$p>\max\{2020,a_1,...,a_{2020}\}$を満たす正整数$n$と素数$p$が存在する. $n$以上の最小の$p-1$の倍数を$m$とする. $p|P(n+(m-n)p)$だから
このとき$0\equiv\displaystyle \sum_{i=1}^{2020}a_i^{n+(m-n)p})\equiv \sum_{i=1}^{2020}a_i^m\equiv\sum_{i=1}^{2020}1=2020\pmod p$. これは$p>2020$に矛盾. よって題意は示された.\qed

多項式$\bmod p$は周期$p$で, 指数関数$\bmod p$は周期$p-1$で変化することはしばしば便利なので多項式と指数関数が絡まる問題においては注意しておこう.(不定方程式とかで無駄な$\bmod$を取ることを回避するときとかにも有効)

\begin{prob}(APMO2021P2)

多項式$P$と正整数$n$に対し$P_n$を$0<a<b\leq n$かつ$n||P(a)|-|P(b)|$を満たす整数の組$(a,b)$の個数とする. 任意の正整数$n$に対して$P_n\leq2021$が成り立つような整数係数多項式$P$を全て求めよ.
\end{prob}

\begin{prob}(ELMO2022P2)\\
以下の条件を満たす最高次係数が$1$の整数係数多項式$P$を全て求めよ:

正整数$a,m$が存在して任意の$n\equiv a\pmod m$なる正整数$n$に対して$$2022\cdot\frac{(n+1)^{n+1}-n^n}{P(n)}$$が整数となる.
\end{prob}

\begin{prob}(Iran 3rd 2012, 2022第3回通信添削P2)

整数係数多項式 $P$ は恒等的には $0$ でないとする. このとき, 以下をみたす素数 $q$
が無限個あることを示せ:
\begin{center}
ある正の整数 $n$ が存在して, $2^n + P(n)$ が $q$ の倍数となる.
\end{center}

\end{prob}

一般に次の定理が知られているらしい. 面白いので載せておく. ($P(n)$とか$2^n+P(n)$は定理の条件を満たす) (prime divisor linear recurrence sequenceとか調べるといっぱい出てくる)
\begin{thm}
$k$を正整数, $u_0,...,u_{k-1},r_1,...r_k$を整数とし$u_n=r_1u_{n-1}+ \cdots + 
r_ku_{n-k}\ (n\geq k)$として$u_n$を定める. もし$u_n$の特性方程式の根$\alpha,\beta\ (\alpha\neq \beta)$であって$\dfrac{\alpha}{\beta}$が$1$の冪乗根であるようなものが存在しないのならば$\mathbb P(u_{\mathbb N})=\mathbb P(\{u_n|n\in \mathbb N\})$は無限集合である.

また, $u_n$の一般項を$u_n=\displaystyle\sum_{i=1}^t f_i(n)\alpha_i^{n}$と表示したとき$t\geq2$となるのならば$\displaystyle\lim_{n\to \infty}\mathbb \max \mathbb P(\{u_n\})=\infty$.
\end{thm}

また素因数の大きさについて, 次のような定理が知られている. (証明はベイカーの対数一形式の理論を用いる)

\begin{thm}
$P(x)$を$0$でない$2$つ以上の根を持つ整数係数多項式とする. このとき$P$に依存するある正定数$c,N$が存在して$N$以上の任意の正整数$n$について$P(n)$は$c\log\log n$以上の素因数を持つ.
\end{thm}

一方多項式については次の未解決問題があるので, これに類似した事実を使おうとする場合は注意が必要である.\\
\textbf{予想.}\ (ブニャコフスキー予想)

$P$を$2$次以上の整数係数既約多項式とする. このとき以下のうち一方が成立するだろう.

($i$)\ $P(n)$が素数となる正整数$n$が無限個存在する.

($ii$)\ ある素数$p$であって任意の正整数$n$について$p|P(n)$が成立するものが存在する.\\

この予想が成立するような多項式の例も成立しない多項式の例も見つかっていない(証明されていない).

一方$P$が一次のときは証明されている. (証明はディリクレの$L$関数という$\zeta$関数の一種を用いる)
\begin{thm}(ディリクレの算術級数定理, ディリクレの素数定理)

$a,b$を互いに素な正整数とするとき$an+b$($n$は正整数)の形で表される素数は無限個存在する. より強く, $x$を実数とし, $\pi_{a,b}(n)$で$an+n$の形で表される素数$p$で$x$以下のものを表すこととするとき$x\rightarrow \infty$のとき
$$\pi_{a,b}(x)\sim \frac1{\varphi(a)}\mathrm{Li}(x)$$
が成立する. ここで$\displaystyle\mathrm{Li}(x)=\int_{2}^x\frac1{\log t}dt$である.(対数積分)
\end{thm}

ちなみに$b=1$のときについては例題$2$で前半の主張を示している. ($(a,b)=(3,2),(4,3)$の場合は前半の主張は簡単に証明できるのでやってみると良いかも知れない)

\begin{expr}(ISL2008N6)

$n^2+1$の素因数として$2n+\sqrt{2n}$より大きいものを持つ正整数$n$が無限に存在することを示せ.
\end{expr}
$n^2+1$が素数となる$n$が無限個存在すれば話は早いがこれは未解決なので諦める必要がある. 同様にある正整数$m$があって$\dfrac{n^2+1}m$が素数となる$n$が無限個存在するか系統の問も上手くいかなさそうという気がする. 発想を転換して素数$p$を固定し, $p|n^2+1$なる最小の$n$を考えてみる. $p>2$の時 $p\equiv1\pmod4$が必要である事に注意.

\textbf{解答}

平方剰余の第一補充則より$p\equiv 1\pmod4$となる素数$p$に対し$p|n^2+1$なる正整数$n$が存在するのでその中で最も小さいものを$n$と置く. $n$は$0<n<p$を満たすとして良く, さらに$p|(p-n)^2+1$に注意すれば$n<\dfrac p2$である. ($n\equiv \pm r\pmod p$の形で$n$の条件が書ける事を使った) よって$p>2n$. もう少し強い評価が必要である. そこで$n=\dfrac{p-1}2-k$と置いて$k$の評価をする. $0\equiv4n^2+4\equiv 4k^2+4k+5 \pmod p$. 従って$4k^2+4k+5=lp$なる正整数$l$が存在する. このとき$p=2n+2k+1=2n+\sqrt{lp-4}$.よって$l\geq3$と出来れば良さそうである.($l=1$ではギリギリ回らない) まず偶奇を考えて$l\neq2$である. $l=1$とすると$\bmod8$より$p\equiv5\pmod 8$が必要. 従って$p\equiv1\pmod8$とすれば$l\geq3$(実際には$l\geq5$)となり, $p$が十分大きい時$\sqrt{lp-4}>\sqrt{2p}>\sqrt{2n}$となる. ディリクレの算術級数定理により$p\equiv1\pmod8$なる素数は無限に存在するので題意は示された.

実際にはこの方法で素因数の評価を$2n+\sqrt{10n}$まで改善出来る.\\

また多項式型だけでなくもっと幅広く使える次の強力な定理も知られている. いくつか定義をする.
\begin{dfn}
$X$を整数の部分集合とし, $a$を整数とするとき$a+X=\{a+x|x\in X\}$とする.
\end{dfn}

\begin{thm}(小林の定理)

$M$を整数からなる無限集合とし$\mathbb P(M)$が有限集合であるとする. また, $a$を$0$でない整数, $m$を$3$以上の整数とし, $b_i$を$M$の元をインデックスに持つ整数とする. $N=\{a+b_t^m\cdot t|t\in M\}$とし$N$が無限集合であるとするとき$\mathbb P(N)$は無限集合である.
\end{thm}

\begin{dfn}
$X$を整数の部分集合とし, $a$を整数とするとき$a+X=\{a+x|x\in X\}$とする.
\end{dfn}

\begin{thm}(小林の定理の系)

$M$を整数からなる無限集合とし, $\mathbb P(M)$が有限集合, $a$を$0$でない整数とする. このとき$\mathbb P(a+M)$は無限集合である.
\end{thm}

ちなみに証明は整数点に関するジーゲルの定理を用いる.
ここまでに出てきた定理の多くは高度すぎて実際に使うのは避けるべきである.(ディリクレの算術級数定理は許容されているとするのが一般的である) 実際に上の定理が使える状況において初等的に如何に議論するかを見る.

\begin{expr}(ISL2012N2)

$P(x)=\displaystyle\prod_{j=1}^9(x+d_j)$とする. ここで$d_j$は相異なる整数である. このときある正整数$N$が存在して, $N$以上の任意の整数$n$について$P(n)$が$20$より大きい素因数を持つことを示せ.

\end{expr}

\begin{expr}(有名事実)

$n,N$を正整数とし, $N$以上の任意の奇素数$p$について$\left(\dfrac np\right)\neq-1$が成立するものとする. このとき$n$は平方数である事を示せ.

\end{expr}

\section{ジグモンディーの定理の利用}
最後にジグモンディーの定理を使った方法を見る.(割とおまけ的な感じ) この定理を使う事で$a^n-b^n$型の数に対して良い条件を満たす素因数の存在をいう事が可能な場合がある. まずはジグモンディーの定理のステートメントを述べる.

\begin{dfn}(原始的素因数)\\
$\{a_n\}$を整数列とする. $a_n$の素因数$p$であって任意の$n$未満の正整数に対して$a_k$が$p$で割り切れないとき$p$を$a_n$の原始的素因数という. また, $a_n$が原始的素因数を持たないとき$\{a_n\}$は$n$-defectiveであると呼ぶ.
\end{dfn}

\begin{thm}(ジグモンディーの定理)\\
$x,y\ (x>y>0)$を互いに素な整数とする.$a_n=x^n-y^n$は以下の例外を除いて原始的素因数を持つ.\\
(1) $n=1,\ x-y=1$のとき($a_1=1$)\\
(2) $n=2$かつ$x+y=2^k$のとき($a_2=2^ka_1$)\\
(3) $n=6,\ x=2,\ y=1$のとき($a_6=63=a_2^2a_3$)\\
また$a_n=x^n+y^n$とすると$n=3,\ x=2,\ y=1$のとき($a_3=9=a_2^2$)を除いて$a_n$は原始的素因数を持つ.
\end{thm}

前半の主張が有名だが後半は$a_{2n}=x^{2n}-y^{2n}$に対して前半の主張を用いれば容易に示せる.

この定理を用いれば瞬殺できる問題を何問か置いておく. (このレベルの問題で答案で使うと多分怒られます)

\begin{expr}(2003初夏P2)\\
$3$以上の奇数$n$と正の整数$x,y$の組であって
$$x^n+y^n=p^m$$をみたす素数$p$と整数$m$が存在するようなものを全て求めよ.
\end{expr}

\textbf{解答}\\
$x,y$が互いに素な場合を考えればよい. $n$が奇数なので$x+y|p^m$. $x+y>1$より$p|x+y$. ジグモンディーの定理により$x^n+y^n$が$p$以外の素因数を持たない条件は$(x,y,n)=(1,1,n),(1,2,3),(2,1,3)$である. あとは$x,y$が公約数を持ちうる事に注意して求める解は$(x,y,n)=(2^t,2^t,n),(3,3^t,2\cdot3^t),(3,2\cdot3^t,3^t)$である.\\

\begin{prob}(JMO2011hoP2)\\
正の整数の組$(a,n,p,q,r)$であって, 等式$$a^n-1=(a^p-1)(a^q-1)(a^r-1)$$を満たすものを全て求めよ.
\end{prob}

\begin{prob}(2000ISL)\\
$(a,m,n)\in \mathbb N^3$であって, $a^m+1|(a+1)^n$を満たすものを全て求めよ.
\end{prob}

一応次のように構図を隠すことも可能かもしれない(?)

\begin{expr}\ \\
$n\leq1000$ かつ以下の $P$ が整数となるような正整数の組 $(n,a_1,a_2,...,a_{10})$はいくつあるか.
$$P=\frac{((2n)!)^{13}}{(a_1^n+1)(a_2^n+1)\cdots(a_{10}^n+1)}$$
\end{expr}
\textbf{解答}

まず次の補題を示す.\\
\textbf{補題}\\
以下の場合を除き, $a^n+1$は $2n+1$ 以上の素因数をもつ：\\
(1)$a=1$\\
(2)$n=1$ かつ $a+1$ が $2$ べき\\
(3)$n=3$ かつ $a=2$\\
\textbf{証明}\\
仮定よりジグモンディの定理から$a^{2n}-1$は$a^i-1$($1\leq i\leq 2n-1$)を割り切らない素因数$p$を持つ. これは$a^n-1$も割り切らないので$a^n+1$の素因数である.もし$p\leq 2n-1$が成り立つと仮定すると, $p$は$a$を割り切らないのでフェルマーの小定理から$a^{p-1}-1$は$p$で割り切れるがこれは$p$が$a^i-1$($1\leq i\leq 2n-1$)を割り切らないことに矛盾する.よって$p\geq 2n+1$となり, 示された.\qed\\

従って$P=\displaystyle\frac{(2n)!^{13}}{(a_1^n+1)(a_2^n+1)...(a_{10}^n+1)}$が整数となるのは,$a_i=1$($i=1,...,10$)のとき, $a_i=2^{k_i}-1$となる正整数$k_i$($i=1,...,10$)が存在し,$n=1$のとき, $a_i=1,2$($i=1,...,10$)かつ$n=3$のときのみあり得る(この$3$条件は排反ではない事に注意).

$(i)\ a_i=1$=1($i=1,...,10$)のとき\\
$n=1,2,...,1000$全てについて$P$は整数となる.よって求める組は$1000$通り.

$(ii)\ a_i=2^{k_i}-1$となる正整数$k_i$($i=1,...,10$)が存在し,$n=1$のとき\\
$P=\displaystyle\frac{2^{13}}{2^{k_1+k_2+...+k_{10}}}$が整数になる条件は$k_1+k_2+...,+k_{10}\leq 13$であるのでこれは${}_{13}\mathrm{C}_{10}=286$通り.

$(iii)\ a_i=1,2$($i=1,...,10$)かつ$n=3$のとき\\
$P$は常に整数となる.よって$2^{10}=1024$通り.\\
$(i),(ii),(iii)$をともに満たすのは$a_i=1$が全ての$i$で成り立ち,$n=1,3$となるときなのでその分を引いて求める解の個数は
$1000+286+1024-2=2308$
である.

\section{もう少し問題を}
最後に練習問題としての置き場を失った問題を何問か置いておこう. なおこの問題の並び順はほぼランダムである.

\begin{prob}(ISL2017N5)\\
$p>q$を満たし,$$\frac{(p+q)^{p+q}(p-q)^{p-q}-1}{(p+q)^{p-q}(p-q)^{p+q}-1}$$が整数となるような素数の組$(p,q)$を全て求めよ.
\end{prob}

\begin{prob}(匿氏の問題)\\
$a,b$を$a>b$なる正整数とする. 任意の正整数$n$に対して$n^a−n^b$が$2^a−2^b$の倍数となるような$(a,b)$を全て求めよ.
\end{prob}

\cite{ジグモンディー}に基づいてジグモンディーの定理の証明の問題を作ってみた.
\begin{prob}(ジグモンディーの定理の証明)\\
$x,y$を$x>y$を満たす互いに素な正整数とし, 数列$a_n$を$a_n=x^n-y^n$で定める. さらに$n$次斉次円分多項式を$\Phi_n(x,y)=y^{\varphi(n)}\phi\left(\dfrac{x}{y}\right)$と定義する. また, $a_n$が原始的素因数を持たないと仮定する.\\
(1)\ $\Phi_n(x,y)$の素因数は$1$または$n$の最大の素因数である事を示せ. \\
(2)\ $n$を正整数とし, $x>1$を実数とする. このとき, 不等式
$(x−1)^{\varphi(n)}\leq\Phi_n(x)<(x+1)^{\varphi(n)}$が成り立つことを示せ. また, 一つ目の不等式における等号は$n=2$でのみ成立することを示せ.\\
(3)\ ジグモンディーの定理を示せ.
\end{prob}

なお, この定理の一般化が知られている.(Y. Bilu, G. Hanrot, P. M. Voutier, Existence of primitive divisors of Lucas and Lehmer numbers, J. Reine Angew. Math. 539 (2001), 75--122.もしくは\url{https://integers.hatenablog.com/entry/2016/05/23/073744}\ を読んでください)次の問題はこの事実が背景にある.

\begin{prob}(OMC40F)\\
$N$ を正整数とします. 数列 $\lbrace a_n\rbrace_{n=1,2,\cdots}$を
$a_1=a_2=1,\ \ a_{n+2}=a_{n+1}+Na_{n}\ \ (n=1,2,\cdots)$で定めるとき, $a_n=p^m$なる合成数 $n$, 素数 $p$, $3$ 以上の整数 $m$ の組の個数を $f(N)$ とします.

このとき
$f(1)+f(2)+\cdots+f\left(\dfrac{3^{12}-1}{2}\right)$
を求めてください. 

ただし, いずれも小数第 $4$ 位で四捨五入した値として, 以下が保証されます.
$$\log_{5}{3}\approx 0.683,\ \ \log_{7}{3}\approx 0.565,\ \ \log_{11}{3}\approx 0.458,\ \ \log_{13}{3}\approx 0.428$$
\end{prob}

\begin{prob}(OMC53F)\\
正の整数の組$(k,n,p,x,y)$は以下の条件をみたします.\\
(1)\ $k$ は $4$ で割って $2$ 余る\\
(2)\ $n\leq 2^{1000}$\\
(3)\ $p$ は $3\leq p\leq 19$なる素数\\
(4)\ $x\leq y$\\
(5)\ $x^k+y^k=p^n$\\
このような組すべてについて $xy$ の総積を考え, これが $2$ で割り切れる回数を $Q$ とします.

$Q$ を $2^{998}-1$で割った余りを $R$ としたとき, $(2^{998}-1)-R$を求めてください.
\end{prob}

次の$5$問は\cite{フィボナッチ平方数}と\cite{平方数の和の素因数}を参考にした.

\begin{prob}(Ion Cucurezeanu)\\
$x^3−x^2+8$が平方数となるような整数$x$を全て求めよ.
\end{prob}

\begin{prob}
$2x^4+2x^2y^2+y^4$が平方数となるような非負整数$x,y$
の組を全て求めよ.
\end{prob}

\begin{prob}
ある立方数より$7$だけ小さいような平方数を全て挙げよ.
\end{prob}

\begin{prob}(フィボナッチ平方数)\\
$F_n,L_n$をそれぞれフィボナッチ数列, リュカ数列とする.(indexを負に拡張して考える)\\
(1)\ 任意の$m,n\in \mathbb Z$に対して, $2|n$かつ$3\not |n$ならば, $F_{m+2n}\equiv−F_m\pmod {L_n}$が成りこと立つを示せ.\\
(2)\ $n\equiv1\pmod4$のとき$F_n$が平方数となるのは$n=1$のときのみであることを示せ.\\
(3)\ $n\equiv3\pmod4$のとき$2F_n$が平方数となるのは$n=3$のときのみであることを示せ.\\
(4)\ $F_n$が平方数となるならば$F_n\in\{0,1,144\}$であることを示せ.
\end{prob}

\begin{prob}(1976 USAMO)\\
方程式$a^2+b^2+c^2=(ab)^2$の整数解を全て決定せよ.
\end{prob}

\begin{prob}(BalkanMO2018P4)\\
$3p^{q-1}+1$が$11^p+17^p$を割り切るような素数の組$(p,q)$を全て求めよ.
\end{prob}

\begin{prob}(ISL1998N5)\\
ある整数$m$が存在して$2^n-1$が$m^2+9$の約数になるような整数$n$を全て求めよ.
\end{prob}

\begin{prob}(IMO1999P4)\\
以下の条件を満たす正整数の組$(n,p)$を全て求めよ.\\
(1)\ $p$は素数\\
(2)\ $n\leq 2p$\\
(3)\ $n^{p-1}|(p-1)^n+1$
\end{prob}

\begin{prob}(2016春P9)\\
$p<q<r<s$である素数の組$(p,q,r,s)$であって,$$p^s-r^q=3$$を満たすものを全て求めよ.
\end{prob}
\textbf{解答}

偶奇を見ることで$p,r$の一方は$2$であり, $p<q<r<s$から$p=2$が分かる.

$p=2$を代入すると$2^s-s^q=3-(s^q-r^q)=3-(s-r)(s^{q-1}+\cdots+r^{q-1})<3-2\cdot2<0$から$s<q\log_2s$が得られる. また, $2^s>2^s-3=r^q$から$s>q\log_2 r$がわかる.

与式を変形すると$2(2^{s-1}-1)=r^q+1$となる. フェルマーの小定理から左辺は$s$で割り切れるので$r^q\equiv-1\pmod s$. 従って$s|r+1$または$s\equiv1\pmod q$となる. $r,s$は$3$以上の素数であり, また$r<s$であるので前者は成立せず, $r\equiv 1\pmod q$と分かる. $s$は奇数である事と合わせて$s=2kq+1$($k$は正整数) と書ける. よって$2(2^{2kq}-1)=r^q+1$. 特に, $2(2^k-1)(2^k+1)|r^q+1$.

$2^k-1<2^k+1<2^{\frac{s}{2q}}+1<\sqrt{s}+1=\dfrac{s}{\sqrt s}+1<\dfrac{s}{\log_2s}+1<q+1$より, $2(2^k-1)(2^k+1)$の素因数$a$は$q$以下である. このとき, $a|r^q+1$は$a|r^{(q,a-1)}+1=r+1$を意味するので$2(2^k-1)(2^k+1)|r+1$である. ($v_q(2(2^k-1)(2^k+1))\leq1$に注意)

$2(2^k-1)(2^k+1)=2^{\frac{s-1}q+1}-2>2^{\frac sq}-2>r-2$より$2(2^k-1)(2^k+1)\geq r-1>\dfrac{r+1}3$ であるので$2(2^k-1)(2^k+1)=r+1$が成立する.

これを$s>q\log_2 r$に代入すると$2kq+1>q\log_2(2^{2k+1}-3)$. よって$1>q\log_2\left(2-\dfrac3{2^{2k}}\right)\geq q\log_2\dfrac54$. これより$q<3.2$つまり$q=3$を得る. $q=3$のとき$s<q\log_2s$は$s<10$を意味し, 従って$(p,q,r,s)=(2,3,5,7)$. これは条件を満たす. 以上より求める解は$(p,q,r,s)=(2,3,5,7)$

\begin{prob}(どちゃ楽数学bot(@solove\_math) Q.175より)\\
$n$を$2$以上の整数とする. 任意の素数$p$に対して$\dfrac{p^n+1}{p+1}$は$n^2$で割り切れない事を証明せよ.
\end{prob}

\begin{prob}(USATSTST2018 P8)\\
$b > 2$を満たす整数$b$で$n^2$が$b^n+1$を割り切るような正整数$n$が無限に存在するようなものを全て求めよ.
\end{prob}

\begin{prob}(2012 USA team selection test P4)

正整数$a,n$であって$a^n-1$を割り切る任意の素数$p$に対し$m<n$である正整数$m$であって$p$が$a^m-1$を割り切るようなものが存在するものを全て求めよ.
\end{prob}

\begin{prob}(USATSTST2008P4)

任意の整数$n$について$n^7+7$は平方数でないことを示せ.
\end{prob}

面白い定理を見つけたので紹介する.
\begin{prob}(Xianzu Lin)

$f:\mathbb N\to\mathbb Z$を次の二条件を満たす関数とする.

\textbf{1.}\ 任意の正整数$a$,非負整数$b$についてある正整数$X$が存在して$f(n)\equiv b\pmod a\Leftrightarrow f(n+X)\equiv b\pmod a$.


\textbf{2.}任意の正整数$a$,非負整数$b$について$n$の関数$f(an+b)$は非有界である.

このとき$\mathbb P(f(\mathbb N))$は無限集合であることを示せ.
\end{prob}

最後のこの問題の初動はこれまでのテクニックを使うが, (想定解では)その後に代数的整数論の高度な結果を使うため解くことを推奨しない. 
\begin{prob}\ \\
以下の条件を満たす整数$n$を全て求めよ:\\
十分大きい全ての整数$x$に対して
$\displaystyle\frac{2(y^4+n)}{x^2+n}$が整数となる整数$y$が存在する.
\end{prob}
\newpage
\section{練習問題解答}(todo:番号を振り直す)
例によって解答がないやつとか適当にリンク貼ってるだけのものとかも多いが許して欲しい.\\
\textbf{2.}\cite{獲得}p.175を見よ.\\
\textbf{4,5.}\cite{合同式に基づく古典的な素数判定法}定理3.1, 3.2を見よ.\\
\textbf{6.}例えば\cite{冪乗}を見よ.\\
\textbf{7.}
\indent 与式を $N$ とおき, $p^2+10$ の適当な素因数を $r$ とする. まず $p\geq 5$ において解が存在しないことを示す.\\
\indent $r\equiv 1\pmod p$ と出来る場合, 簡単な計算によって $p\geq 23$ である. $r=kp+1$ とおけば以下は整数である.
$$\frac{p^2+10}{r}=\frac{p(p-10k)}{kp+1}+10$$
ここで $p$ と $kp+1$ は互いに素であるから, 特に $p-10k$ は $kp+1$ で割り切れる. ここで $p$ が素数であることより $p-10k\neq 0$ であるから $p+10k\geq |p-10k|\geq kp+1$ が必要であるが, これは $(k-1)(p-10)\leq 11$ と変形されるため不適である. したがって, 以下すべての $r$ について $r\not\equiv1\pmod p$ であるとしてよい.\\
\indent このとき, さらに $r\not\equiv 1\pmod q$ と出来る場合, Fermatの小定理より $5^{r-1}\equiv 1\equiv 5^{3pq} \pmod r$ が従う.
ここで $\gcd(r-1,3pq)=\gcd(r-1,3)$ であるから, $r$ は $5^3-1=124$ の約数であり, $p^2+10$ が奇数であることと合わせて $r=31$ である. しかし, 平方剰余の相互法則より $p^2\equiv -10\pmod{31}$ とはなり得ないことが確認できるから, 以下すべての $r$ について $r\equiv 1 \pmod q$ であるとしてよい.\\
\indent このとき, $p^2+10\equiv 1\pmod q$ である. 明らかに $q\neq 3$ であるから, $-1$ は ${\rm mod}\ q$ において平方剰余である. ここで $q\equiv 1\pmod 4$ のとき, $3^q-8$ が $5$ で割り切れるため不適であるから, $q=2$ である. このとき上と同様に $r$
は $5^6-1$ の約数であり, 特に $r\not\equiv 2\pmod 3$ である. すなわち $p^2+10\not\equiv 2\pmod 3$ であるが, これは $p\geq 5$ に反する. したがって, 以上より $p\geq 5$ において解が存在しないことが示された.\\
\indent 続いて, $p=2,3$ として $q\geq 5$ において解が存在しないことを示す. Fermatの小定理より $3^q-8\equiv -5\neq 1 \pmod q$ であるから, $3^q-8$ の素因数 $s$ であって $s\not\equiv 1\pmod q$ なるものが存在する. このとき, 上と同様に $r$ は $5^6-1$ または $5^9-1$ の約数, すなわち $r=3,7,19,31,829$ である. $r=3$ は明らかに不適. $r=7,19,31$ は $3^q\equiv 8\pmod r$ のとき常に $q$ は $3$ で割り切れるから不適である. $r=829$ のとき, 平方剰余の相互法則より $3^q$ は ${\rm mod}\ r$ において平方剰余であるが, 一方で第二補充則より $8$ は平方非剰余であるから矛盾する.\\
\indent よって, 以上より $p,q\leq 3$ であり, この範囲で確かめれば求める解は $(p,q)=(2,2),(2,3),(3,2)$ である.\\
\textbf{8.}(公式解説でも良いが)例えば\cite{獲得}p.172または\cite{精選}p.189を見よ.\\
\textbf{9.}\cite{冪乗}または\cite{精選}p.138を見よ.\\
\textbf{10.}\cite{冪乗}または\cite{精選}p.190を見よ.\\
\textbf{12.}概略だけ. $p$が奇素数のときを考える. prime-constraintで$n$の素因数が決まり, $n$の偶奇でそれぞれ$\dfrac{n^2+1}p, \dfrac{n+1}p$が整数だと分かる. ここで$n$の約数のうち素因数として$p$で割って$1$余るもののみもつものの最大のものを考えて, 頑張ると出来る(はず).\\
\textbf{14.}まず以下の補題を示す.\\
\textbf{補題.}　$5$ 以上の奇数 $n$ に対し, $2^n+37^n$は $4$ で割って $3$ 余る $3$ 以外の素因数をもつ.\\
\textbf{証明.}　$n$ が $9$ で割り切れるとき, $2^9+37^9\equiv 1\pmod{4}$であり, LTEの補題より $2^9+37^9$は $3$ で $3$ 回割り切れることと併せて成立する. したがって, 以下 $n$ は素因数 $p\geq 5$ をもつとする. このとき $\textrm{mod}\ 9$ を考えれば $2^p+37^p$は $3$ でちょうど $1$ 回割り切れるから, 同様に $\textrm{mod}\ 4$を考えて所望の結論の成立が分かる.\qed\\

$b$ が奇数のとき, 両辺を $\textrm{mod} ~ 8$ で見て$ a=1$ を得る. このとき $39-3^b=c^2$より $(a,b,c)=(1,1,6)$ である.

以下 $b$ が偶数であるとする. $a$ が $5$ 以上の奇数のとき, 補題より $2^a+37^a$ は $4$ で割って $3$ 余る $3$ 以外の素因数 $p$ をもつが, $-1$は $\textrm{mod}~p$で平方非剰余であるからこれは矛盾である. したがって $a$ が奇数ならば $a=1,3$ である.

$a=1$ は上で確かめた. $a=3$ のとき, $\textrm{mod} ~ 9$ より $b\geq 2$ で $c$ は $3$ の倍数であることに留意すれば
$5629-3^{b-2}=\left(\dfrac{c}{3}\right)^2$
$b=2,4,6,8$ をそれぞれ調べることで, $(a,b,c)=(3,8,210)$を解として得る.

$a$ が偶数のとき, $\textrm{mod} ~ 3$ より$b=0$ がわかる. $a\geq2$ のとき
$37^a< 2^a+37^a-1< (37^{a/2}+1)^2$
より不適であり, $a=0$ のとき $(a,b,c)=(0,0,1)$を解として得る.

以上より求める解は$(a,b,c)=(1,1,6),(3,8,210),(0,0,1)$である.\\
\textbf{18.}\cite{kkkaaa}を見よ.\\
\textbf{19.}\cite{ジグモンディー}を見よ.\\
\textbf{20.}\url{https://onlinemathcontest.com/contests/omc040/editorial/277}\\
\textbf{21.}\url{https://onlinemathcontest.com/contests/omc053/editorial/216}なおここでは高度な結果の利用を避け, 愚直に解いている.\\
\textbf{25.}\cite{フィボナッチ平方数}を見よ.\\
\textbf{33.}$\mathbb Q(i),\ \mathbb Q(i,\sqrt n),\ \mathbb(i,\sqrt[4]n)$での素イデアルの完全分解の条件と完全分解する素イデアルの集合で体の拡大が決まることに留意すると$n=-k^4,\ 4k^4\ (k\in \mathbb Z)$のみが条件を満たすことが分かる.
四乗剰余の相互法則とディリクレの算術級数定理を用いてもよい.

\begin{thebibliography}{99}
\bibitem{フィボナッチ平方数}{Fibonacci 平方数が 0, 1, 144 のみであることの初等的証明}(\url{https://mathlog.info/articles/447})
\bibitem{平方数の和の素因数}{特殊な形の不定方程式で有用な解の限定法}(\url{https://mathlog.info/articles/664})
\bibitem{冪乗}{冪乗が出てくる整数問題}(\url{https://mathlog.info/articles/2941})
\bibitem{kkkaaa}{$n^a-n^b$が$2^a-2^b$の倍数となる$a,b$を求める}(\url{https://mathlog.info/articles/2624})
\bibitem{合同式に基づく古典的な素数判定法}合同式に基づく古典的な素数判定法(\url{https://math.jp/wiki/合同式に基づく古典的な素数判定法})
\bibitem{ジグモンディー}{Zsigmondyの定理}(\url{https://integers.hatenablog.com/entry/2016/12/30/183841})
\bibitem{30}{30：Lucas数列と原始的約数}(\url{https://integers.hatenablog.com/entry/2016/05/23/073744})
\bibitem{すごい論文}Y. Bilu, G. Hanrot, P. M. Voutier, Existence of primitive divisors of Lucas and Lehmer numbers, J. Reine Angew. Math. 539 (2001), 75--122.
\bibitem{獲得}獲得金メダル！国際数学オリンピック
\bibitem{精選}数論の精選104問
\bibitem{eurosky}Eurosky 杯(\url{https://twitter.com/sugaku_2357/status/1493510732142751745})
\bibitem{dio}Diophantine Approximation
(\url{https://www.uni-due.de/~adg350u/Skripte/DioApp.pdf})
\bibitem{kobayashi}On Existence of Infinitely Many Prime Divisors in a Given Set\\(\url{https://projecteuclid.org/euclid.tjm/1270215162})
\bibitem{lrpd}On the greatest prime factor of terms of a linear recurrence sequence\\(\url{https://x.gd/r5MCU})
\bibitem{top}Xianzu Lin\ \ Prime Divisors Of Sequences Of Integers
\end{thebibliography}
\end{document}